\documentclass[12pt]{article}
\usepackage[a4paper, total={6.5in, 9.5in}]{geometry}
\usepackage{graphicx}			
\usepackage{amssymb}
\usepackage{amsmath}			
\usepackage{fancyhdr}
\usepackage{enumerate}       
\usepackage{dirtytalk} 
\usepackage[english]{babel}
\usepackage{amsthm}
\usepackage{xcolor}
\usepackage[shortlabels]{enumitem}
\usepackage{titlesec}
\titleformat{\section}{\normalfont\Large\bfseries}{\thesection}{0pt}{\Large}
\setlength{\parindent}{0pt}
\setlength{\parskip}{0.8em}

\newtheorem{theorem}{Theorem}[section]
\newtheorem{corollary}{Corollary}[theorem]
\newtheorem{lemma}[theorem]{Lemma}
\newtheorem{definition}[theorem]{Definition}


\begin{document}
% \renewcommand{\thesection}{Exercise \arabic{section}}


\setcounter{section}{0}
\renewcommand{\thesection}{}


\parindent = 0pt


      {\bf Analysis 1 HW 6 Part 1}
%==================

\section{5.2}

$f:(a, b)\to \mathbb{R}$.\\
$f'(x)>0$ on $(a, b)$.\\
$g\equiv f^{-1}$.\\

{\bf To prove:} $f$ is strictly increasing in $(a, b)$. 

\begin{proof}
      

Let $x_{1}, x_{2}\in(a, b)$ such that $x_{1}<x_{2}$. We know that $f$ is differentiable on $[x_{1}, x_{2}]$. The mean value theorem guarantees the existence of $p\in(x_{1}, x_{2})$ satisfying
$$
\frac{{f(x_{2})-f(x_{1})}}{x_{2}-x_{1}}=f'(p).
$$
Since $f'(p)> 0$, we have $f(x_{2})-f(x_{1})> 0$. Thus, $f$ is strictly increasing $(a, b)$.
\end{proof}
{\bf To prove:} $g$ is differentiable.

\begin{proof}
Let $I$ be the image of $f$ on $(a, b)$. Let $y\in I$. We have to show that
$$
\lim_{ t \to y } \frac{{g(y)-g(t)}}{y-t} = \frac{1}{f'(g(y))}.
$$
Let $x\equiv g(y)$, and $t'\equiv g(t)$. The expression inside the limit then becomes
$$
\frac{{x-t'}}{f(x)-f(t')}.
$$
We know that
$$
\lim_{ t' \to x } \frac{{x-t'}}{f(x)-f(t')}= \frac{1}{f'(x)},
$$

since $f'(x)>0$. Let $\epsilon>0$. There exists $\delta_{1}>0$ such that 
$$
t'\in(x-\delta_{1}, x+\delta_{1}) \implies \left\lvert  \frac{{x-t'}}{f(x)-f(t')}- \frac{1}{f'(x)} \right\rvert <\epsilon.
$$
By taking a suitable closed interval $[x-\zeta, x+\zeta]$, we get that $g$ is continuous at $y$, since continuous bijective maps on compact sets are homeomorphisms. So, 
$$
\lim_{ t \to y } g(t)=g(y).
$$
It follows that there exists $\delta_{2}>0$ such that 
$$
t\in(y-\delta_{2}, y+\delta_{2})\implies t'\in(x-\delta_{1}, x+\delta_{1}).
$$
\end{proof}





\pagebreak

\section{5.8}

{\bf To prove:} If $f'$ is continuous on $[a, b]$, $f$ is uniformly differentiable on $[a, b]$.


\begin{proof} 


WLOG, let $a\leq t < x \leq b$. Since $f'$ is uniformly continuous on $[a, b]$, there exists $\delta>0$ such that $0<x-t<\delta \implies |f'(t)-f'(x)|<\epsilon$. Now, from the mean value theorem, there exists $t<t'<x$ such that
$$
\frac{{f(t)-f(x)}}{t-x}=f'(t').
$$
So, $0<x-t<\delta \implies 0<x-t'<\delta$ $\implies$ $|f'(t')-f'(x)|<\epsilon$.
\end{proof}

{\bf To prove:} If $f'$ is continuous on $[a, b]$, $f$ is strongly differentiable on $[a, b]$.
\begin{proof} 
Let $\epsilon>0$. Consider a sequence $(x_{n}, y_{n})\to(t, t)$ in $\mathbb{R}^{2}$. Since $f'$ being continuous on $[a, b]$ implies $f$ is uniformly differentiable on $[a, b]$, there exists $\delta>0$ such that
$$
|x_{n}-y_{n}|<\delta \implies \left|\frac{{f(x_{n})-f(y_{n})}}{x_{n}-y_{n}}-f'(y_{n})\right|< \frac{\epsilon}{2}.
$$
From slot-wise convergence, we have $(x_{n})\to t$ and $(y_{n})\to t$. Thus, there exists $N_{1}, N_{2}\in \mathbb{Z}^{+}$ such that $n>N_{1}\implies |x_{n} - t|< \delta/2$ and $n>N_{2} \implies |y_{n}-t|< \delta/2$. From the triangle inequality, we have
$$
n>\max\{N_{1}, N_{2}\} \implies|x_{n}-y_{n}|\leq |x_{n}-t|+|y_{n}-t|<\delta.
$$
Now, since $f'$ is continuous, there exists $\delta_{2}>0$ such that $|y-t|<\delta_{2}\implies |f'(y)-f'(t)|< \epsilon/2$. Also, there exists $N_{3}$ such that $n>N_{3} \implies |y_{n}-t|<\delta_{2}$. Let $N=\max\{ N_{1}, N_{2} , N_{3}\}$. Then, we have
$$
n>N \implies \left|\frac{{f(x_{n})-f(y_{n})}}{x_{n}-y_{n}}-f'(t)\right| \leq \left|\frac{{f(x_{n})-f(y_{n})}}{x_{n}-y_{n}}-f'(y_{n})\right|+|f'(y_{n})-f'(t)|<\epsilon.
$$
Thus, 
$$
\lim_{ (x, y) \to (t, t) } \frac{{f(x_{n})-f(y_{n})}}{x_{n}-y_{n}}=f'(t).
$$
\end{proof}
\pagebreak




\section{5.17}

$f$ is a real, thrice differentiable function on $[-1, 1]$.\\
$f(-1)=0$, $f(0)=0$, $f(1)=1$, $f'(0)=0$.\\

{\bf To prove:} $f'''(x)\geq 3$ for some $x\in(-1, 1)$.

\begin{proof}
On applying Taylor's theorem with $\alpha=0$ and $\beta=1$, we get

\begin{align*}
f(1)=\frac{f''(0)}{2}+\frac{f'''(t)}{6} \\
\implies 6=3f''(0)+f'''(t).
\end{align*}

for some $t\in(0, 1)$. With $\beta=-1$, we get

\begin{align*}
f(-1)=\frac{f''(0)}{2}-\frac{f'''(s)}{6} \\
\implies f'''(s)=3f''(0).
\end{align*}

for some $s\in(-1, 0)$. Thus, we have
$$
f'''(s)+f'''(t)=6.
$$
This is impossible if $f'''(x)<3$ for all $x\in(-1, 1)$. 
\end{proof}

\end{document}